\chapter*{はじめに} \addcontentsline{toc}{chapter}{はじめに}

本書は，著者が中心となって開発しているプログラミング言語Egisonを通して提唱しているパターンマッチ指向プログラミングというプログラミング・パラダイムをできるだけコンパクトに紹介することを目的として執筆された．
パターンマッチまわりのEgisonの構文と，それらの構文を使ったプログラミングテクニックを紹介する．

Egisonはより直感的にアルゴリズムを表現したいという動機で発案されたプログラミング言語のアイデアを実証するためにつくられたproof of conceptのプログラミング言語である．
Egisonに実装されているこれらのアイデアのうち，もっとも重要な機能は本書の主題であるパターンマッチである．
Egisonのパターンマッチの特徴は，ユーザーにより拡張可能でかつ，非線形パターン（パターン変数に束縛した値を同じパターン内で参照するパターン）をバックトラッキングにより効率的に処理することである．
本書の目的は，この機能がいかに表現力が豊かで，応用範囲がひろいのか読者に著者と同じくらいに感じてもらうことである．
%本書ではふれないが，他の重要な機能には，微分幾何の計算を表現するのに便利なテンソルの添字記法をプログラム中で使うための機能がある．

本書内のサンプルプログラムは，現在のEgisonの文法を使って記述されている．
今後の開発でEgisonの文法が変わり続ける可能性は高い．
しかし，サンプルプログラムの裏のアイデアの根幹は普遍的なものであるので，安心して読んでほしい．

本書がきっかけとなって，プログラミング言語，とくにパターンマッチの研究をする研究仲間が読者の中から現れてくれたらとてもうれしい．

\begin{flushright}
  2020年1月6日
  
  江木 聡志
\end{flushright}
